\documentclass[leqno]{article}
\input{../../../../_shared/pack_preamble.tex}

\djsQbtHeader{Mark Scheme Fixture}
\djsSitePath{Core Pure}{Polar Coordinates}

\djsLastUpdated{18 October 2026}
\begin{document}
\djsFrontMatter
\begin{enumerate}[
  label=\textbf{\arabic*.},
  leftmargin=*,
  topsep=0pt,
  partopsep=0pt,
  parsep=0pt
]

% ---- Question 1 ----
\questionitem
% [Source: _QBT__Fixture, Q1]

\begin{enumerate}[label=\textbf{(\alph*)}]
  \item On the same diagram, sketch the curve $C_1$ with polar equation
  \[
  r=2(1+\cos\theta),\quad 0\leq\theta\leq2\pi
  \]
  and the curve $C_2$ with polar equation $\theta=\frac{2\pi}{3}$. \marks{3} \\
  \item Find the area of the smaller region bounded by $C_1$ and $C_2$. \marks{4} \\
\end{enumerate}

\newpage

% ---- Question 2 ----
\questionitem
% [Source: _QBT__Fixture, Q2]

One end of a light elastic string, of natural length $2a$ and modulus of elasticity $2mg$, is attached to a fixed point $O$. The other end of the string is attached to a particle $P$ of mass $m$. The string is vertical with $P$ directly above $O$. The particle $P$ is pulled vertically upwards until the extension of the string is $2a$. The particle $P$ is then released from rest. \\
\begin{enumerate}[label=\textbf{(\alph*)}]
  \item Find the speed of $P$ when it is at a distance $\dfrac{3}{2}a$ above $O$. \marks{4} \\
  \item Find the initial acceleration of $P$ when it is released from rest. \marks{2} \\
\end{enumerate}

\newpage

% ---- Question 3 ----
\questionitem
% [Source: _QBT__Fixture, Q3]

A plane passes through the points $C(1,-1,2)$, $D(3,0,1)$ and $E(2,2,2)$. \\

The point $A$ has coordinates $(4,-2,7)$ and the point $B$ is the reflection of $A$ in the plane. \\

Find the coordinates of the point $B$. \marks{8}

\newpage

% ---- Question 4 ----
\questionitem
% [Source: _QBT__Fixture, Q4]

Two smooth uniform spheres $P$ and $Q$ have equal radii and masses $0.4\,\mathrm{kg}$ and $0.6\,\mathrm{kg}$ respectively. They move on a smooth horizontal surface and collide obliquely. \\

Immediately before the collision, their velocities are
\[
\mathbf{u}_P=(6\mathbf{i}+2\mathbf{j})\,\mathrm{m\,s^{-1}},\qquad \mathbf{u}_Q=(-2\mathbf{i}+2\mathbf{j})\,\mathrm{m\,s^{-1}}
\] \\
where $\mathbf{i}$ and $\mathbf{j}$ are perpendicular horizontal unit vectors. \\

At the instant of impact, the vector from the centre of $P$ to the centre of $Q$ is a positive multiple of $\mathbf{i}+\mathbf{j}$. The total kinetic energy of the two spheres decreases by $2.88\,\mathrm{J}$ during the collision. \\
\begin{enumerate}[label=\textbf{(\alph*)}]
\item Find
\begin{enumerate}[label=(\roman*)]
\item the velocity of $P$ immediately after the collision,
\item the velocity of $Q$ immediately after the collision,
\item the coefficient of restitution between the spheres.
\end{enumerate}
Justify your answers, explaining which velocity components are unchanged by the collision. \marks{16} \\
\item Find the angle through which the direction of motion of $P$ is deflected by the collision. Give your answer in degrees to one decimal place. \marks{2} \\
\end{enumerate}

\end{enumerate}
\end{document}
